\section{Related Work}

\label{sec:related}

\subsection{Satellite-Based Change Detection}

Change detection from satellite imagery has evolved through three generations.
First-generation methods used algebraic operations on spectral bands: image differencing, image ratioing, and change vector analysis \citep{singh1989review}.
Second-generation methods employed machine learning classifiers (random forests, SVMs) on hand-engineered features \citep{zhu2017change}.
Third-generation methods use deep learning for end-to-end change detection.

Among deep learning approaches, \citet{daudt2018fully} introduced fully convolutional Siamese networks for change detection.
\citet{chen2020spatial} proposed a spatial-temporal attention network for bi-temporal change detection.
\citet{zheng2021change} demonstrated that transformer architectures achieve state-of-the-art change detection accuracy.
However, most work focuses on urban change detection rather than ecological monitoring, and processes image pairs rather than temporal sequences.

\subsection{Deforestation Monitoring}

Global Forest Watch \citep{hansen2013high} provides the most widely used deforestation monitoring, based on Landsat-derived tree cover change maps updated annually.
GLAD alerts \citep{hansen2016humid} improved timeliness to monthly updates using Landsat-8 data.
Planet-NICFI \citep{planet2020nicfi} provides high-resolution basemaps for tropical forest monitoring but relies on visual interpretation rather than automated detection.

\citet{maretto2020spatio} applied deep learning to Sentinel-2 for Amazon deforestation detection, achieving 88\% accuracy.
\citet{ortega2021near} developed near-real-time deforestation detection using Sentinel-1 SAR data, robust to cloud cover but at lower spatial resolution.
None of these systems integrate change detection with ecological impact assessment or conservation-prioritized alerting.

\subsection{Landscape Ecology Metrics}

Landscape ecology quantifies habitat pattern through metrics such as patch area, edge density, connectivity indices, and fragmentation measures \citep{mcgarigal2012fragstats}.
\citet{haddad2015habitat} demonstrated that habitat fragmentation reduces biodiversity by 20--75\% depending on scale and ecosystem type.
\citet{fahrig2017ecological} analyzed the relative effects of habitat loss versus fragmentation per se, finding that loss is the dominant driver but fragmentation contributes additional negative effects through edge effects and isolation.

Remote sensing-based fragmentation assessment has been limited to static analyses \citep{vogt2007mapping} rather than continuous monitoring.
Our HFI provides the first real-time fragmentation metric derived from change detection.

\subsection{Conservation Alert Systems}

Early warning systems for conservation include the FIRMS fire detection system \citep{giglio2016active}, which detects active fires from MODIS and VIIRS data within 3 hours.
\citet{finer2018combating} developed the Real-Time Deforestation Alert System (RAISG) for the Amazon using Sentinel-1 radar data.
\citet{brancalion2020global} proposed a prioritization framework for forest restoration but not for active threat response.

% ============================================================
% 3. Temporal Siamese Network
% ============================================================
